%%%%%%%%%%%%%%%%%
% This is an sample CV template created using altacv.cls
% (v1.7.2, 28 August 2024) written by LianTze Lim (liantze@gmail.com). Compiles with pdfLaTeX, XeLaTeX and LuaLaTeX.
%
%% It may be distributed and/or modified under the
%% conditions of the LaTeX Project Public License, either version 1.3
%% of this license or (at your option) any later version.
%% The latest version of this license is in
%%    http://www.latex-project.org/lppl.txt
%% and version 1.3 or later is part of all distributions of LaTeX
%% version 2003/12/01 or later.
%%%%%%%%%%%%%%%%

%% Use the "normalphoto" option if you want a normal photo instead of cropped to a circle
% \documentclass[10pt,a4paper,normalphoto]{altacv}

\documentclass[10pt,a4paper,ragged2e,withhyper]{altacv}
%% AltaCV uses the fontawesome5 and simpleicons packages.
%% See http://texdoc.net/pkg/fontawesome5 and http://texdoc.net/pkg/simpleicons for full list of symbols.

% Change the page layout if you need to
\geometry{left=1.25cm,right=1.25cm,top=1.5cm,bottom=1.5cm,columnsep=1.2cm}

% The paracol package lets you typeset columns of text in parallel
% \usepackage{paracol}

% Change the font if you want to, depending on whether
% you're using pdflatex or xelatex/lualatex
% WHEN COMPILING WITH XELATEX PLEASE USE
% xelatex -shell-escape -output-driver="xdvipdfmx -z 0" sample.tex
\iftutex 
  % If using xelatex or lualatex:
  \setmainfont{Roboto Slab}
  \setsansfont{Lato}
  \renewcommand{\familydefault}{\sfdefault}
\else
  % If using pdflatex:
  \usepackage[rm]{roboto}
  \usepackage[defaultsans]{lato}
  % \usepackage{sourcesanspro}
  \renewcommand{\familydefault}{\sfdefault}
\fi

% Change the colours if you want to
\definecolor{SlateGrey}{HTML}{2E2E2E}
\definecolor{LightGrey}{HTML}{666666}
\definecolor{DarkPastelRed}{HTML}{450808}
\definecolor{PastelRed}{HTML}{8F0D0D}
\definecolor{GoldenEarth}{HTML}{E7D192}
\colorlet{name}{black}
\colorlet{tagline}{PastelRed}
\colorlet{heading}{DarkPastelRed}
\colorlet{headingrule}{GoldenEarth}
\colorlet{subheading}{PastelRed}
\colorlet{accent}{PastelRed}
\colorlet{emphasis}{SlateGrey}
\colorlet{body}{LightGrey}

% Change some fonts, if necessary
\renewcommand{\namefont}{\Huge\rmfamily\bfseries}
\renewcommand{\personalinfofont}{\footnotesize}
\renewcommand{\cvsectionfont}{\LARGE\rmfamily\bfseries}
\renewcommand{\cvsubsectionfont}{\large\bfseries}


% Change the bullets for itemize and rating marker
% for \cvskill if you want to
\renewcommand{\cvItemMarker}{{\small\textbullet}}
\renewcommand{\cvRatingMarker}{\faCircle}
% ...and the markers for the date/location for \cvevent
% \renewcommand{\cvDateMarker}{\faCalendar*[regular]}
% \renewcommand{\cvLocationMarker}{\faMapMarker*}


% If your CV/résumé is in a language other than English,
% then you probably want to change these so that when you
% copy-paste from the PDF or run pdftotext, the location
% and date marker icons for \cvevent will paste as correct
% translations. For example Spanish:
% \renewcommand{\locationname}{Ubicación}
% \renewcommand{\datename}{Fecha}


%% Use (and optionally edit if necessary) this .tex if you
%% want to use an author-year reference style like APA(6)
%% for your publication list
% \input{pubs-authoryear.tex}

%% Use (and optionally edit if necessary) this .tex if you
%% want an originally numerical reference style like IEEE
%% for your publication list
\input{pubs-num.tex}

%% sample.bib contains your publications
\addbibresource{sample.bib}
% \usepackage{academicons}\let\faOrcid\aiOrcid
\begin{document}
\name{Md Adil}
\tagline{Self-taught Software Engineer | Architecting scalable systems and cloud-native platforms.}
%% You can add multiple photos on the left or right
% \photoR{2.8cm}{Globe_High}
% \photoL{2.5cm}{Yacht_High,Suitcase_High}

\personalinfo{%
  % Not all of these are required!
  \email{adil.sudo@gmail.com}
  \phone{+91 98 1812 8797}
  \location{Pune, India}
  \homepage{md-adil.github.io}
  % \twitter{@twitterhandle}
  \xtwitter{@_MdAdil}
  \linkedin{md-adil}
  \github{md-adil}
  
  %% You can add your own arbitrary detail with
  %% \printinfo{symbol}{detail}[optional hyperlink prefix]
  % \printinfo{\faPaw}{Hey ho!}[https://example.com/]

  %% Or you can declare your own field with
  %% \NewInfoFiled{fieldname}{symbol}[optional hyperlink prefix] and use it:
  % \NewInfoField{gitlab}{\faGitlab}[https://gitlab.com/]
  % \gitlab{your_id}
  %%
  %% For services and platforms like Mastodon where there isn't a
  %% straightforward relation between the user ID/nickname and the hyperlink,
  %% you can use \printinfo directly e.g.
  % \printinfo{\faMastodon}{@username@instace}[https://instance.url/@username]
  %% But if you absolutely want to create new dedicated info fields for
  %% such platforms, then use \NewInfoField* with a star:
  % \NewInfoField*{mastodon}{\faMastodon}
  %% then you can use \mastodon, with TWO arguments where the 2nd argument is
  %% the full hyperlink.
  % \mastodon{@username@instance}{https://instance.url/@username}
}

\makecvheader
%% Depending on your tastes, you may want to make fonts of itemize environments slightly smaller
% \AtBeginEnvironment{itemize}{\small}

%% Set the left/right column width ratio to 6:4.
\columnratio{0.6}

% Start a 2-column paracol. Both the left and right columns will automatically
% break across pages if things get too long.
% \begin{paracol}{2}
\cvsection{Experience}

\cvevent{Sr. Software Architect}{Bajaj Finserv Health}{Jan 2020 -- Ongoing}{Pune}
\begin{itemize}
\item Leading squad of Smart Reports and OCR with Agentic AI.
\item An in-house DLS called \textbf{capsule} (\href{https://design.bajajfinservhealth.in}{design.bajajfinservhealth.in})
\item Introduced \textbf{Micro-frontend} and \textbf{ArgoCD}
\item Handled security concerns, including addressing DDoS attacks and
SQL vulnerabilities.
\item Implemented deployment optimizations resulting in an \textbf{80\%} reduction
in deployment time and a \textbf{99\%} decrease in ACR size. These enhancements have led to significant cost
savings and notable improvements in efficiency.
\item Establishing coding guidelines and ensuring adherence to them
throughout the organization.
\end{itemize}

\begin{itemize}
\item Introduced Next.js \& re-architected the frontend
\item Designed and developed reusable React hooks and components that
have gained significant adoption across the organization.
\item Took ownership of designing and developing B2C and Labs,
starting from the initial concept and building them entirely from
scratch.
\end{itemize}
\divider


\cvevent{Tech Lead \& Full Stack}{BigRadar}{April 2017  -- Oct 2020}{Delhi}
\begin{itemize}

\item Built a full-fledged application from scratch independently, striving to
achieve industry standards and optimal performance despite limited
resources.
\item Developed a customized ORM solution for MongoDB and optimized a
wrapper for Express.js, enhancing the efficiency and functionality of
both technologies.
\end{itemize}

\divider

\cvevent{Sr. Software Engineer}{Edunuts}{Aug 2015 -- Apr 2017}{Delhi}

Overseeing and enhancing the company's technology stack through
the strategic selection of appropriate technologies and frameworks.

\medskip

Tech Stack
\cvtag{PHP}
\cvtag{Laravel}
\cvtag{Slim}
\cvtag{MySQL}
\cvtag{Node.js}
\cvtag{MongoDB}

\divider


\newpage


\textbf{Other Companies:}

\begin{itemize}
    \item \href{https://apkaakhbar.in/}{Apka Akhbar}
    \item Bigly Technologies
    \item GuruInfoways
\end{itemize}

\cvsection{Case Studies}
% --- Case Study 1 ---
\cvevent{PHP-FPM → Laravel Octane Migration}{EdTech Platform}{}{}

\textbf{Tech Stack:} Laravel Octane, Swoole, Redis, MySQL, PHP

\medskip

\textbf{Challenge:}
Slow responses (800ms avg), poor concurrency, and high CPU usage under PHP-FPM.

\medskip

\textbf{Solution:}
Migrated to Laravel Octane + Swoole with coroutine execution, process-level caching, and DB connection pooling.

\medskip

\textbf{Outcomes:}


\begin{itemize}
  \item Response times improved from 800ms → 40–60ms
  \item 5× higher concurrency handling
  \item 30–40\% infra cost reduction
\end{itemize}
\divider
% --- Case Study 2 ---
\cvevent{Monolith → Microservices Migration}{EdTech \& SaaS Platforms}{}{}

\textbf{Tech Stack:} Kubernetes, Docker, Node.js, Go, Laravel, Jenkins

\medskip

\textbf{Challenge:}
Monolith causing scaling issues, slow deployments, and feature delivery delays.

\medskip

\textbf{Solution:}
Implemented strangler-fig migration to microservices with full observability on Kubernetes.

\medskip

\textbf{Outcomes:}
\begin{itemize}
  \item Improved fault isolation across services
  \item Clear ownership and autonomous deployments
  \item Independent database + service-level scaling
\end{itemize}
\divider
% --- Case Study 3 ---
\cvevent{Real-time Observability Platform}{Clients}{}{}

\textbf{Tech Stack:} Grafana, ELK Stack, Loki, Tempo, Prometheus, Kubernetes

\medskip

\textbf{Challenge:}
Poor visibility in distributed systems; MTTR over 4 hours for critical issues.

\medskip

\textbf{Solution:}
Built end-to-end observability with tracing, centralized logging, metrics, and Grafana dashboards.

\medskip

\textbf{Outcomes:}

\begin{itemize}
  \item MTTR reduced from 4+ hours → under 30 minutes
  \item Achieved 99.99\% availability
  \item Enabled proactive issue detection before user impact
\end{itemize}



\cvsection{Projects}

\cvevent{Yet Another Chatting Software}{\href{https://yacs.in}{yacs.in}}{May 13, 2019}{Hometown}

Talk to anyone without sharing contacts

\medskip

Tech stack \cvtag{Typescript} \cvtag{React} \cvtag{Node.js} \cvtag{MongoDB} \cvtag{Redis} \cvtag{Protobuff}

\divider
\cvevent{Minimajs}{\href{https://minima-js.github.io/}{minima-js.github.io}}{Sep, 2023}{Pune}
 MinimaJS is a cutting-edge Node.js framework designed to empower developers to construct contemporary web applications with exceptional efficiency and elegance.
 
\medskip

\begin{itemize}
    \item Functional Programming
    \item TypeScript Foundation
    \item Modern JavaScript Adoption
    \item Direct Request Access
\end{itemize}

\divider

\cvevent{jiq (Javascript inline query)}{\href{https://md-adil.github.io/jiq/}{md-adil.github.io/jiq}}{Nov, 2020}{Pune}
Use existing javascript knowledge to query or modify almost any type data, files and directory.
\medskip


\medskip

% \cvsection{A Day of My Life}

% % Adapted from @Jake's answer from http://tex.stackexchange.com/a/82729/226
% % \wheelchart{outer radius}{inner radius}{
% % comma-separated list of value/text width/color/detail}
\wheelchart{1.5cm}{0.5cm}{%
  6/8em/accent!30/{Sleep,\\beautiful sleep},
  3/8em/accent!40/Learn,
  8/8em/accent!60/Daytime job,
  2/10em/accent/Gym and relaxation,
  5/6em/accent!20/Spending time with family
}

% use ONLY \newpage if you want to force a page break for
% ONLY the current column
% \newpage
%% Switch to the right column. This will now automatically move to the second
%% page if the content is too long.
% \switchcolumn

\cvsection{My Life Philosophy}

With 11+ years across backend engineering, cloud infrastructure, performance optimization, and developer tooling, I specialize in designing and implementing end-to-end solutions that scale.


\cvsection{Most Proud of}

\cvachievement{\faTrophy}{Design Language System (\href{https://design.bajajfinservhealth.in/}{Capsule})}{Architected and development of Capsule — a scalable design system used across Bajaj Finserv Health products, achieved 99\% lighthouse score.}

\cvachievement{\faTrophy}{Architecture}{
Built scalable cloud-native architectures (1M+ users), led microservices modernization, created full K8s + GitOps platforms, solved end-to-end engineering problems, and built widely-used open-source tools.

}

\cvsection{Strengths}

\begin{itemize}
    \item Self-Taught \& Passionate
    \item Open Source Contributor \& Creative Technologist
    \item Tooling \& Automation
    \item Technical Leadership
    \item Security-First Thinking
    \item DevOps \& Deployment Optimization
    \item Deep understanding of scalable system architecture.
    \item Quick Learner - I can learn anything.
\end{itemize}

\cvsection{Technology}

\textbf{Frontend}

\smallskip

\cvtag{TypeScript}
\cvtag{JavaScript}
\cvtag{React}
\cvtag{Next.js}
\cvtag{Vue.js}
\cvtag{Socket.io}
\cvtag{WebRTC}

\divider

\textbf{Backend}

\smallskip

\cvtag{Node.js}
\cvtag{Express.js}
\cvtag{Nest.js}
\cvtag{Objection.js}
\cvtag{Mongoose}
\cvtag{PHP}
\cvtag{Laravel}
\cvtag{Slim}

\divider

\textbf{Styling \& UI}

\smallskip

\cvtag{CSS}
\cvtag{SASS}
\cvtag{Tailwind}
\cvtag{Bootstrap}

\divider
\newpage
\textbf{DevOps / Cloud}

\smallskip

\cvtag{Docker}
\cvtag{Kubernetes}
\cvtag{AWS}
\cvtag{Azure}
\cvtag{GCP}
\cvtag{Nginx}
\cvtag{Azure Pipeline}
\cvtag{Github Actions}

\divider

\textbf{Databases}

\smallskip

\cvtag{MySQL}
\cvtag{PostgreSQL}
\cvtag{MongoDB}
\cvtag{Redis}

\divider

\textbf{Messaging / Streaming}

\smallskip

\cvtag{Kafka}
\cvtag{RabbitMQ}
\cvtag{RTMP}
\cvtag{HLS}

\divider

\textbf{Observability}

\smallskip

\cvtag{Grafana}
\cvtag{Prometheus}
\cvtag{ELK Stack}
\cvtag{Loki}
\cvtag{Tempo}


\divider


Least Experience 
\medskip

\cvtag{C++}
\cvtag{Embedded systems programming}
\cvtag{React Native}
\cvtag{Arduino}
\cvtag{Python}
\cvtag{Electron}

\cvsection{Tools}
\cvtag{Git} \cvtag{Webpack} \cvtag{Rollup} \cvtag{Vite}

\newpage

\cvsection{Education}

\cvevent{SOFTWARE ENGINEERING}{School of Internet}{2010 -- Still Learning}{}

\divider

\cvevent{BACHELOR OF COMMERCE}{University of Delhi}{2010 -- June 2013}{}

\cvsection{Open source}

I enjoy building tools that solve real developer problems — this mindset naturally led me to open-source.
Notable projects include:

\medskip

\begin{itemize}
    \item \href{https://minima-js.github.io/}{Minimajs}
    \item \href{https://www.npmjs.com/package/react-read-otp}{react-read-otp} (\~1k weekly downloads)
    \item \href{https://www.npmjs.com/package/ebx}{ebx}
\end{itemize}

\smallskip

\href{https://github.com}{Explore 150+ on Github}

\smallskip

Over \textbf{15 Packages} on \href{https://www.npmjs.com/~adil.sudo}{Npmjs}



\cvsection{Hobbies \& Interests}

\begin{itemize}
    \item \textbf{Programming as a Hobby} – I enjoy experimenting with new technologies,
building side projects, and exploring open-source contributions beyond
work.
    \item \textbf{Movies \& Storytelling} – I appreciate well-crafted narratives, which also
inspires my interest in user experience and product design.
    \item \textbf{Logic Puzzles & Brain Teasers} – Keeps my problem-solving sharp and
helps me think outside the box in debugging and systems design.
\end{itemize}
% \end{paracol}


\end{document}
